\section{The Traveler — Ways and Journeys}

\subsection{Domain and Epithets}

The Traveler is the eternal guide of the road, guardian of those who walk the paths between what is and what might be. The Traveler is not merely one who shows the way—they are the way, existing in the pause between steps and in the choice of which path to take when roads fork. Every journey is both physical and spiritual; to move from one place to another is to transform, and the road itself becomes a teacher.

The Traveler has no temples, only crossroads and waystations. No priests, only guides and couriers. No scripture, only the accumulated wisdom of a thousand roads walked and remembered.

\subsection{Cultural Origins}

The Traveler is invoked wherever journeys begin and end:

\begin{itemize}
\item In \textbf{Kahfagia}, the Traveler is the Lantern-Bearer, whose light guides ships through uncharted waters. Pilots light a beacon before every voyage.
\item In \textbf{Ykrul}, the Traveler is the Hoof-Print, the spirit of the steppe road that appears as a single arrowhead left at a crossroads.
\item In \textbf{Aeler}, the Traveler is the Ways-Stone, carved into every mile marker by True-Masons to ensure the road remembers those who walk it.
\item In \textbf{Fhara}, the Traveler is the Oasis Father, the guide who never drinks—but always finds water.
\end{itemize}

\subsection{Warrior Archetype: The Road-Warden}

A Runic Warrior of the Traveler is not a sentinel or a judge. They are a \textit{Road-Warden}—a warrior who carries the road in their feet and their steel. Their Codex is a waystone or a compass. Their Thiasos is a spirit-horse that never tires, a road-dog that knows every shortcut, or a sentient blade that has walked more miles than its wielder.

\subsection{Codex Form}

The Road-Warden's Codex is a tool of navigation and passage:

\begin{itemize}
\item \textbf{Waystone:} A small carved stone, warm to the touch. It always points to the nearest road or safe passage.
\item \textbf{Compass:} A brass compass whose needle never points north—it points to the nearest threshold. Useful for finding doors, gates, and borders.
\item \textbf{Road-Staff:} A walking staff carved with route maps. Tapping it on the ground reveals the safest path forward.
\end{itemize}

\subsection{Thiasos Form}

\begin{itemize}
\item \textbf{Spirit-Horse:} A horse that never tires and leaves no tracks. It can gallop for a full day without rest, and it knows the way home.
\item \textbf{Road-Dog:} A grey hound that can smell the difference between a safe road and an ambush. It whines when danger is near.
\item \textbf{The Waystone Compass (Sentient):} A compass that speaks in the sound of footsteps. It can find a shortcut (reduce travel clock by 2 segments once per journey), but it will ask for a story in return.
\end{itemize}

\subsection{Core Mechanic: The Road's Memory}

When the Road-Warden travels a road they have walked before, they gain +1 die to all navigation, survival, and stealth rolls on that road. The road remembers them.

When the Road-Warden helps another traveler find their way (offering directions, an escort, or a map), they gain 1 Boon. This Boon can only be used on rolls involving travel, navigation, or escape.

\subsection{Sample Rites}

\subsubsection{Low Rite: Rite of Road-Sense (4 XP)}

\textbf{Tags:} [NAVIGATION], [AWARENESS], [TRAVEL]

\textbf{Threshold:} A road-nail or waystone pebble.

Unerringly pick the fastest safe route in Near or Far. Gain +1 die to avoid ambushes or delays. Create a 4-segment Path Memory clock to ignore difficult terrain once.

\textbf{Obligation:} 1

\subsubsection{Low Rite: Rite of the Traveler's Boon (5 XP)}

\textbf{Tags:} [MOVEMENT], [BOND], [HASTE], [TRAVEL]

\textbf{Threshold:} Thread tied at the wrist.

Ignore one level of difficult terrain or bureaucracy; +1 Effect to travel or escape checks. If shared, create a 2-segment Shared Journey bond with an ally.

\textbf{Obligation:} 1

\subsubsection{Standard Rite: Rite of the Waymark (8 XP)}

\textbf{Tags:} [MARK], [GUIDANCE], [PATH], [TRAVEL]

\textbf{Threshold:} A chalk mark or small cairn.

Declare a lane as permitted or easy. Allies gain improved Position/Effect or ignore one obstacle. Create a 6-segment Marked Path clock.

\textbf{Obligation:} 2

\subsubsection{Standard Rite: Rite of the Bridge Between (7 XP)}

\textbf{Tags:} [TELEPORT], [PATH], [BOND], [TRAVEL]

\textbf{Threshold:} Two pinches of road-dust clapped together.

Relocate a willing target within Far along a visible or named route. Unwilling targets may resist (Body + Resolve DV 3). Create a 4-segment Pathway Established clock.

\textbf{Obligation:} 2

\subsubsection{High Rite: Rite of the Wayfarer's Covenant (14 XP)}

\textbf{Tags:} [VOW], [COMMAND], [BOND], [FEAR], [TRAVEL]

\textbf{Threshold:} Waystones from multiple regions.

Bind present parties to safe passage. While honored: +2 Effect on travel, reroll one failed travel roll per scene. Breaking the oath inflicts Harm 1 (Fatigue) and marks the breaker as \textit{Oathbreaker of the Road} (-2 dice to travel rolls).

\textbf{Obligation:} 3

\subsection{Patron-Specific Talents}

\subsubsection{Shortcut (4 XP)}

\textbf{Requirement:} The Traveler

Once per session when moving from one location to another (including during travel montages), declare a hidden shortcut that reduces the remaining travel clock by 2 segments. The GM may describe a minor cost or complication (e.g., "the path is steep—mark 1 Fatigue").

\subsubsection{Road Memory (6 XP)}

\textbf{Requirement:} The Traveler

You never forget a road you have walked. You can draw a perfect map of any route you have traveled, including shortcuts, hazards, and waystations. This is a passive ability.

\subsubsection{Wayfinder's Blessing (8 XP)}

\textbf{Requirement:} The Traveler, Tier II+

Once per journey, you may declare that a road you are traveling is blessed. For the duration of the journey, you and your party suffer -1 DV to all navigation and survival rolls, and the first ambush is automatically detected. The GM gains 2 SB (the road remembers the favor).

\subsection{Roleplaying a Road-Warden}

\begin{itemize}
\item You do not stay in one place. The road is your home; the horizon is your destination.
\item You may be called to guide caravans, escort refugees, or find lost travelers. Your knowledge of the road is worth more than gold.
\item You may struggle with the cost of movement. Every journey leaves something behind. The road remembers what you lost.
\item You may be asked to close a road—to deny passage, to seal a gate, to turn someone away. The Traveler does not forbid stillness, but stillness is not the way.
\end{itemize}

\subsection{Sample NPC: Road-Warden Cerran of the Broken Trail}

\textbf{Archetype:} Road-Warden

\textbf{Patron:} The Traveler

\textbf{Codex:} A compass that points to thresholds, not north

\textbf{Thiasos:} A grey hound that smells ambushes

\textbf{Personality:} Quiet, watchful, restless. He has not slept in a bed for three years. He does not intend to start.

\textbf{Conflict:} A road he once guided is now haunted. Something follows travelers. Cerran cannot remember the route—the road has forgotten him. He needs to walk it again to remind the stones who he is.

\subsection{GM Guidance: Intrusions for The Traveler}

\begin{itemize}
\item A road the Road-Warden has walked before now leads somewhere else—a different building, a different forest, a different year.
\item The Road-Warden's boots wear out instantly. They must walk barefoot for a scene, feeling every stone.
\item A stranger approaches on the road and asks for directions. If the Road-Warden lies, they become lost for 1d4 hours. If they tell the truth, the stranger follows.
\item The Road-Warden's compass spins without stopping. The threshold is close—but it does not want to be found.
\end{itemize}

The thread held. The water carried. That is not a Rite. That is grace.
